
\sgaStableTarget{sga6.e10.sec.7.13}\subsection*{7.13. Relations with @@M0@@-adic cohomology}
Recall (SGA 5 VII) that, for every locally free Module of finite type @@M1@@ on a scheme @@M2@@, and every prime number distinct from the residue characteristics of @@M3@@, one defines @@M4@@-adic Chern classes (for @@M5@@)
@@M6@@
satisfying an additivity formula
@@M7@@
when @@M8@@ is an extension of a locally free module @@M9@@ by a locally free module @@M10@@. From this additivity one concludes that the total Chern class @@M11@@ arises from a homomorphism
\sgaStableTarget{sga6.e10.eq.7-13-0}@@M12@@
(where the target is a group under multiplication, as explained in V 1, and where @@M13@@ is the ``\emph{naive}'' @@M14@@). One obtains canonical maps
\sgaStableTarget{sga6.e10.eq.7-13-1}@@M15@@

Now let @@M16@@ be a \emph{proper} flat scheme over @@M17@@, with \emph{regular fibers}, and suppose that @@M18@@ is \emph{henselian}. Assume that @@M19@@ the residue characteristic of @@M20@@. One can then form the diagram
\sgaStableTarget{sga6.e10.eq.7-13-2}@@M21@@
where the horizontal arrows are the ``Chern class'' maps for @@M22@@ and @@M23@@ respectively, and where @@M24@@ is the composite homomorphism
@@M25@@
